% ============================================================================
%  3. RESULTS
%
%  Esta es la sección donde normalmente viven las figuras y las tablas.
%  Referencia rápida de los comandos de aapaper.sty; el detalle está en
%  la skill global aa-paper (references/figures.md y references/tables.md).
%
%  ------------------------------------------------------------------ FIGURAS
%  A&A admite tres anchos y solo tres:
%
%    \aafig[label]{archivo}{caption}[colocación]        88 mm, una columna
%    \aafigside[label]{archivo}{caption}[colocación]   120 mm, caption al lado
%    \aafigwide[label]{archivo}{caption}[colocación]   180 mm, dos columnas
%    \aafigmulti[label]{f1,f2,f3}{caption}[colocación] varios gráficos
%
%  El archivo se busca en figures/ (ver \graphicspath en main.tex).
%  Los floats a dos columnas hay que declararlos TEMPRANO en el fuente:
%  LaTeX solo puede empujarlos hacia adelante, nunca hacia atrás.
%
%  Requisitos de las imágenes:
%    - eps, pdf, jpg, png o tiff. Bitmaps a 250-300 dpi mínimo, no 92.
%    - Recortar el espacio vacío alrededor.
%    - Sin grids ni fondos de color salvo que aporten información.
%    - Nada de rojo+verde juntos (accesibilidad para daltonismo).
%    - Sin líneas finas; en minúsculas el texto dentro de la figura.
%    - Los símbolos se explican en el caption, no dentro de la figura.
%
%  ------------------------------------------------------------------- TABLAS
%    \aatable[label]{archivo.csv}{caption}[nota]        una columna
%    \aatablewide[label]{archivo.csv}{caption}[nota]    dos columnas
%    \aatablelong[n]{archivo.csv}{caption}{label}       multipágina
%    \aatablebib{(1)~\citet{...}; (2) \citet{...}.}     refs. bajo la tabla
%
%  El CSV va en tables/. Primera fila = encabezado. Las celdas admiten LaTeX
%  crudo ($M_\odot$, \tablefootmark{a}). Si un valor lleva coma, comillas
%  dobles alrededor.
%
%  Los comandos ya aplican las reglas A&A: caption arriba, \hline\hline /
%  \hline / \hline (NO booktabs), columnas a la izquierda, sin líneas
%  verticales. La tabla completa no debe exceder 23.5 cm de alto.
%
%  Notas al pie de tabla:
%    - general:     el 4.o argumento de \aatable pasa a \tablefoot{}
%    - referenciada: \tablefootmark{a} en la celda del CSV, y en la nota
%                    \tablefoottext{a}{texto de la nota}
% ============================================================================

\section{Results}
\label{sec:results}

Texto de la sección. Ejemplo vivo de los comandos: la \figref{fig:demo} y la
\tabref{tab:demo} de abajo compilan y se expanden con \texttt{tools/flatten.py}.
Bórralas cuando empieces a escribir de verdad.

\aafig[fig:demo]{demo}{Figura de demostración.}

\aatable[tab:demo]{tables/demo.csv}
  {Tabla de demostración generada desde \texttt{tables/demo.csv}.}
  [Nota general de la tabla.
   \tablefoottext{a}{Nota referenciada desde una celda del CSV.}]

% ----------------------------------------------------------------------------
%  Ejemplos — descomentar y adaptar
% ----------------------------------------------------------------------------
%
% \aafig[fig:orbit]{orbit.pdf}{Projected orbit of \object{S2} around
%   \object{Sgr A*}. Filled circles mark the astrometric epochs.}
%
% \aafigwide[fig:residuals]{residuals.pdf}{Residuals of the fit in right
%   ascension (\emph{left}) and declination (\emph{right}).}
%
% \aafigside[fig:posterior]{posterior.pdf}{Marginalised posterior distribution
%   of the black hole mass.}
%
% \aatable[tab:elements]{tables/elements.csv}
%   {Best-fitting orbital elements.}
%   [Uncertainties are the 68\% credible intervals.
%    \tablefoottext{a}{Held fixed during the fit.}]
%
% \aatablewide[tab:comparison]{tables/comparison.csv}
%   {Comparison with previous determinations.}
%
% \aatablelong[1]{tables/catalogue.csv}{Full astrometric catalogue.}{tab:catalogue}
% ----------------------------------------------------------------------------
